\documentclass[11pt,a4paper]{article}
\usepackage[T1]{fontenc}
\usepackage[utf8]{inputenc}
\usepackage[margin=25mm]{geometry}
\usepackage{amsmath,amssymb,mathtools}
\usepackage{mathrsfs}
\usepackage{booktabs,longtable,array}
\usepackage{xcolor,listings}
\usepackage[colorlinks=true,linkcolor=blue!50!black,urlcolor=blue!50!black,citecolor=blue!50!black]{hyperref}
\allowdisplaybreaks[2]
\setlength{\emergencystretch}{2em}
\numberwithin{equation}{section}
\newcommand{\QQ}{\mathbb Q}
\newcommand{\CC}{\mathbb C}
\newcommand{\dd}{\mathrm d}
\newcommand{\bdel}{\bar\partial_{\bar\psi}}
\newcommand{\tht}{\vartheta}
\newcommand{\del}{\partial}
\newcommand{\Pg}[2]{P_{#1}^{(#2)}}
\newcommand{\wt}{\operatorname{wt}}
\newcommand{\an}{\mathrm{an}}
\newcommand{\hol}{\mathrm{hol}}
\newcommand{\code}[1]{\texttt{\detokenize{#1}}}
\lstset{basicstyle=\ttfamily\small,breaklines=true,columns=fullflexible,
  keepspaces=true,frame=single,rulecolor=\color{black!25},showstringspaces=false}
\title{An Exact BCOV Algebra in Yamaguchi--Yau Variables\\
\large Definitions, conventions, signs, normalization factors, and implementation audit}
\author{Companion to \texttt{bcov\_yy.py}, version 0.2.0 (frame audit)}
\date{19 October 2026}
\begin{document}
\maketitle

\begin{abstract}
This note specifies the exact mathematical problem solved by the accompanying
quintic BCOV implementation. It derives the differential ring in
$u,v_1,v_2,v_3,X$, the normalized covariant derivative, and the three polynomial
anomaly equations, accounting explicitly for the minus sign and factor
$1/2$ in the recursion. The genus-two calculation is expanded coefficient by
coefficient. Holomorphic inputs, algebraic consequences, and conventions for
integration constants are distinguished throughout. The note also documents
the sparse exact representation, shared arithmetic recipes, validation, and
limitations of the implementation. It does not claim that unfixed boundary
data have been determined by the anomaly equation.
\end{abstract}

\tableofcontents
\newpage

% This is an included, supplementary section in bcov_yy_conventions_and_derivations_v2.tex.
\section{Convention ledger and exact cross-frame audit (19 October 2026)}
\label{sec:ledger2026}
This section records the normalization and gauge reconciliation among original
Yamaguchi--Yau (YY), Huang--Klemm--Quackenbush (HKQ), and the small-radius
coordinate used in the Gu/Klemm free-energy source files. The original YY
BCOV recursion below remains our \emph{canonical computational convention}.
Coordinates and gauges are metadata, not implicitly interchangeable symbols.

\subsection{Coordinates, Yukawa tensors, and normalized amplitudes}
Write $s$ for the original YY fifth-root coordinate, $\psi=s^5$ for HKQ, and
$z=(3125\psi)^{-1}$ for the small large-radius coordinate. In particular,
\begin{equation}\label{eq:2026coordinates}
 \theta_s=5\theta_\psi=-5\theta_z,\qquad
 X_{\rm YY}=X_{\rm HKQ}=\frac{1}{1-\psi},\qquad
 X_z=\frac{1}{1-3125z},\qquad X_{\rm HKQ}=1-X_z.
\end{equation}
The holomorphic Yukawa coupling $C_{ijk}$ is a cubic tensor: it transforms with
\emph{three} coordinate Jacobians. Set $C_t=t^3C_{ttt}$ for each coordinate $t$,
using a \emph{common} choice of holomorphic three-form. Then
\begin{equation}
 C_s=125 C_\psi,\qquad C_\psi=-C_z,
 \qquad C_z=5X_z,\quad C_{zzz}=5X_z/z^3,
 \quad P_g^{(t)}=C_t^{g-1}F_g^{(t)}.
\end{equation}
A further constant rescaling of the holomorphic three-form changes the
individual Yukawa and section-valued free energy but cancels in the normalized
$P_g$, provided it is implemented consistently. The $\kappa=5$ A-model frame
and factors of $(2\pi i)^3$ should not be silently combined.

Under the metric pullback $G_{\psi\bar\psi}=|dz/d\psi|^2G_{z\bar z}$
(with the same K\"ahler frame),
\begin{equation}
 A_\psi=-A_z-2,\quad B_{p,\psi}=(-1)^p B_{p,z},
 \quad A_s=5A_\psi+4,\quad B_{p,s}=5^p B_{p,\psi}.
\end{equation}
The polynomial generators consequently satisfy
\begin{align}\label{eq:2026gens}
(u,v_1,v_2,v_3,X)_\psi
 &=(-u,-v_1,v_2,-v_3,1-X)_z,\\
(u,v_1,v_2,v_3,X)_s
 &=(5u_\psi,5v_{1,\psi},25v_{2,\psi},125v_{3,\psi},X_\psi).
\end{align}
If $T$ acts by substituting the variables in the right-hand side, the corresponding
free-energy polynomial transformations are
\begin{align}\label{eq:2026polys}
P_g^{\rm HKQ}(v_1,v_2,v_3,X)
 &=(-1)^{g-1}P_g^{z}(-v_1,v_2,-v_3,1-X),\\
P_g^{\rm YY}(v_1,v_2,v_3,X)
 &=5^{3g-3}P_g^{\rm HKQ}(v_1/5,v_2/25,v_3/125,X).
\end{align}
These are involutive/rational polynomial frame maps, checked symbolically at
$g=2$ and with multiple independent polynomial round trips. They include
the overall Yukawa sign; changing only $X$ is insufficient.

\subsection{Propagator gauge: supplied five holomorphic functions}
Set $\mu=1/3125$, $\kappa=5$ and use $X_z$ in the following equations.
The holomorphic functions are precisely the supplied Gu/Klemm choice:
\begin{equation}\label{eq:2026gauge}
 q^z_{zz}=-\frac8{5z},\quad q_z^{zz}=\frac z{25},\quad
 q_z^z=-\frac6{625},\quad q_{zz}=\frac2{25z^2},\quad
 q_z=\frac2{3125z}.
\end{equation}
In this gauge the three modified propagators are
\begin{align}\label{eq:2026prop}
 S^{zz}&=-\frac{z^2}{5X_z}\left(v_{1,z}+\frac35\right),\\
 \widetilde S^{z}&=\frac{z}{5X_z}\left(v_{2,z}+\frac2{25}\right),\\
 \widetilde S&=\frac{1}{10X_z}
 \left[v_{3,z}+(1-X_z)\left(v_{2,z}+\frac4{125}\right)\right].
\end{align}
The two derivative identities determining $\widetilde S^z,\widetilde S$ have
been independently reduced to zero using exact rational SymPy expressions.
The fifth $q_z$ enters the larger BCOV differential system; it is
\emph{recorded but not tested} by those two propagator identities.
The original background notes had a repeated $\widetilde S^z$ label on the
left-hand side of the third propagator equation. It should read
$\widetilde S$. Likewise $z\partial_z=X(X-1)\partial_X$ only
\emph{when acting on functions of $X$ alone}.

\subsection{Genus-two anchor and period/mirror-map normalization}
The exact polynomial from the source $\code{ff_2.txt}$ after the above coordinate
transformation is the HKQ expression
\begin{align}\label{eq:2026P2}
 P_2^{\rm HKQ}={}&\frac{25}{144}-\frac{625}{288}v_1
 +\frac{25}{24}v_1^2-\frac5{24}v_1^3
 -\frac{625}{36}v_2+\frac{25}{6}v_1v_2+\frac{350}{9}v_3\notag\\
 &-\frac{5759}{3600}X-\frac{167}{720}v_1X
 +\frac16v_1^2X-\frac{475}{12}v_2X
 +\frac{41}{3600}X^2-\frac{13}{288}v_1X^2+\frac{X^3}{240}.
\end{align}
Substitution into the original YY frame reproduces the published $P_2$
and the program's $\code{reference_genus2()}$.
For the Frobenius periods with $z\to0$,
\begin{align}
\omega_0(z)&=\sum_{n\ge0}\frac{(5n)!}{(n!)^5}z^n,\\
\sigma_1(z)&=\sum_{n\ge1}\frac{(5n)!}{(n!)^5}
 5(H_{5n}-H_n)z^n,\\
 t&=\frac{\omega_1}{\omega_0},\quad
 \omega_1=\frac{\omega_0\log z+\sigma_1}{2\pi i},\quad
 q=e^{2\pi i t}=z\exp(\sigma_1/\omega_0).
\end{align}
In particular $z(q)=q-770q^2+171525q^3-81623000q^4+\cdots$.
The correctly normalized A-model limit is
\begin{equation}\label{eq:2026hol}
 F_g^A=\omega_0^{2g-2}
 \left(\frac{1-\psi}{5\psi}\right)^{g-1}
 P_g^{\rm HKQ}\big|_{\mathrm{hol}},\qquad \psi=(3125z)^{-1}.
\end{equation}
Our older Mathematica \code{BCOVpropagators.txt} helper used a fifth-root
variable in its $C_{fff}$ definition but substituted the HKQ variable
$1/(3125z)$ directly for it. That substitution omits a fifth root and,
for $C_{sss}$, a cubic Jacobian. Its object called $\code{Cf[z]}$ was
\emph{an inverse Yukawa} in that helper, whereas the source $\code{ff_g.txt}$
uses $C_{zzz}=5X_z/z^3$. These two quantities must never share an
unqualified symbol in executable code.

\subsection{Exact enumerative check: genera 2--15, degrees 1--6}
For $g\ge2$ the instanton contribution is extracted by
\begin{equation}\label{eq:2026gv}
 [q^d] F_g^A
 =\sum_{m\mid d}m^{2g-3}
  \sum_{h=0}^{g} n^h_{d/m}\,
  [y^{g-h}]\left(\frac{2\sin(\sqrt y/2)}{\sqrt y}\right)^{2h-2}.
\end{equation}
For $h=0$ the displayed coefficient prescription is
$[y^g](2\sin(\sqrt y/2)/\sqrt y)^{-2}$; it accounts for the
$\lambda^{-2}$ genus-zero contribution correctly. Genus-one GV data
contribute no higher-genus coefficient because the sine power is zero.
The source is \code{quinticGVDTPT.m} and its \code{GVTab}, not a manually typed
subset. Running \code{verify_physical_qseries.py} confirms \textbf{all 84
positive-degree coefficients}, $g=2,\dots,15$ and $d=1,\dots,6$, with
exact rational zero residuals. For example, $[q]F_2^A=2875/240=575/48$:
the free-energy coefficients are rational genus-$g$ Gromov--Witten
quantities, not the integer GV invariants themselves.

The 61-by-44 $\code{GVTab}$ and the partial 113-by-31
$\code{GVTabExt}$ disagree at 29 numerical overlaps, beginning at
$(g,d)=(26,16)$. These conflicts have \emph{not} been adjudicated;
our free-energy sources stop at genus 15. Do not treat an unknown entry
in the extended table as zero, nor mistake its sparse genus-64 row for a
complete genus-64 enumeration.

\subsection{Bernoulli and arithmetic invariants after frame correction}
Define $P_g=P_g^{\rm an}+a_g(X)$ in the original YY frame,
with $P_g^{\rm an}(0,0,0,X)=0$, and $a_g(X)=\sum_n a_{g,n}X^n$.
The two independently re-tested formulas are
\begin{align}\label{eq:2026bernoulli}
 b_1&=\frac16,\quad
 b_g=\frac{1}{2g+1}\sum_{h=1}^{g-1}{2g\choose 2h}b_hb_{g-h},
 \qquad b_g=|B_{2g}|,\\
 a_{g,0}&=100\,625^{g-1}\frac{b_gb_{g-1}}{2g(2g-2)(2g-2)!}.
\end{align}
For $g\ge3$, put
\begin{equation}
 S_g=\sum_{h=2}^{g-2}(2h-2)(2g-2h-2)a_{h,0}a_{g-h,0},\qquad
 I_g=\frac{23}{25}\frac{(-625)^{g-1}b_g}{2g(2g-2)!},
\end{equation}
and then
\begin{align}\label{eq:2026a1}
 a_{g,1}={}&-\frac{577}{625}(g-1)a_{g,0}
 +(2g-4)\left(62-\frac{24g}{5}\right)a_{g-1,0}
 -\frac{12}{5}S_g-I_g.
\end{align}
Genus two supplies the separate seed $a_{2,1}=-28795/144$.
We extracted $[X^0],\ldots,[X^4]$ from all 14 source amplitudes
without full multivariate expansion. Equations \eqref{eq:2026bernoulli} and
\eqref{eq:2026a1} pass \textbf{28 of 28 independent rational equalities}
for $g=2,\ldots,15$.

The original YY and HKQ pure coefficients are related by
$a_{g,n}^{\rm YY}=5^{3g-3}a_{g,n}^{\rm HKQ}$, so ratios
$a_{g,n}/a_{g,0}$ agree in these frames. By contrast, under $X\mapsto1-X$
the \emph{individual} coefficients mix by an integral, upper-triangular
binomial matrix with determinant $\pm1$. Therefore the least common
multiple of their reduced denominators is exactly invariant:
\begin{equation}\label{eq:2026lcd}
 \mathop{\rm lcm}_n\operatorname{den}(a_{g,n}^{z})
 =\mathop{\rm lcm}_n\operatorname{den}(a_{g,n}^{\rm HKQ}).
\end{equation}
Between YY and HKQ, the non-$5$-primary part of this denominator lattice
is unchanged. The assertion follows algebraically from the integral
inverse matrix, and was checked against full pure-$X$ polynomials for
genera 2--8. This protects denominator-lattice patterns, not individual
face-value sequences, from the coordinate inversion. General propagator
gauge transformations may mix the \emph{primitive/ambiguity split} and
must be handled separately.

\subsection{Provenance, regression suite, and unsolved problems}
The v0.2 archive includes the unchanged YY recursion and the additional
frame maps, tests, and source-derived data. The immutable
\code{FramedPolynomial} wrapper carries the genus and frame and rejects
additions of expressions with incompatible frame metadata. The principal modules are:
\begin{quote}\small\raggedright
\code{bcov_yy.py}, \code{quintic_frames.py},
\code{test_frames.py}, \code{test_gu_gauge.py},
\code{audit_coefficients.py}, \code{audit_denominators.py},
\code{verify_physical_qseries.py}.
\end{quote}
The identity checks include $P_2$ in both frames, the two Gu gauge
derivative equations, formal round trips, Bernoulli coefficients to genus
15, full pure-$X$ denominator lattices to genus 8, and the 84 independent
$q$-series checks. The source free energies and source GV table both come
from the same research provenance; the comparison is algebraically independent
but is not a proof of mathematical correctness of all the tabulated invariants.

We have \emph{not} determined holomorphic ambiguities without external
boundary inputs, settled the 29 conflicting high-genus integers, proved a
new small-integer uniqueness principle, or certified an all-genus Bernoulli
formula for the middle powers $n\sim g$. The independently established
fixed-low-$X$ jet recursions through genus 53 belong to the original YY
normalization and remain valid there; the present audit independently
rechecks the source free energies through genus 15.


\section{Scope and the distinction between inputs and consequences}

The implementation specializes to the mirror quintic and to the original
$\psi$ convention of Yamaguchi and Yau (YY) \cite{YY}. The equations below
refer to that convention unless a different coordinate or gauge is explicitly
introduced. The four central objects must not be conflated:
\begin{align}
F_g(\psi,\bar\psi)&:
  \text{the B-model amplitude in a chosen Hodge-line gauge},\\
P_g&=C^{g-1}F_g:
  \text{the normalized polynomial amplitude},\\
P_g^{\an}&:
  \text{the implementation's particular polynomial primitive},\\
F_g^A(t)&:
  \text{the holomorphic A-model amplitude after the frame conversion}.
\end{align}
Here $C=\psi^3 C_{\psi\psi\psi}$, not $C_{\psi\psi\psi}$ itself.

The geometric inputs are the quintic Picard--Fuchs operator, special geometry,
the choice of holomorphic terms that closes the differential ring, the
genus-one one-point function, and any supplied higher-genus boundary data.
The program derives higher anomaly polynomials from those inputs by exact
polynomial arithmetic. It neither proves the special-geometry identities
from the construction of the Calabi--Yau family nor independently solves all
physical boundary conditions.

In particular, the four pure-$X$ coefficients supplied at genus two are
reference boundary data from the complete YY polynomial. Their presence
must not be described as a prediction of the anomaly recursion alone.
The note makes the origin of each such coefficient explicit.

\section{Coordinates, discriminants, and Euler derivatives}
\label{sec:coordinate}

The mirror family is obtained from the quintic polynomial
\begin{equation}
p(x,\psi)=x_1^5+x_2^5+x_3^5+x_4^5+x_5^5
          -5\psi x_1x_2x_3x_4x_5
\end{equation}
with its usual mirror quotient. Set
\begin{equation}
z=(5\psi)^{-5}=\frac{1}{5^5\psi^5},\qquad
\tht=\psi\del_\psi,\qquad \theta_z=z\del_z.
\end{equation}
The factor $5^5=3125$ is part of the definition of $z$.
Since
\begin{equation}
\del_\psi z=-\frac{5z}{\psi},\qquad
\tht z=-5z,
\end{equation}
the chain rule gives
\begin{equation}\boxed{\tht=-5\theta_z.}\label{eq:theta-coordinate}\end{equation}
This accounts for the division by $-5$ in \code{YYRing.theta_z}.
It implements $z\del_z$, not $\del_z$. To obtain $\del_z$ one would
also need multiplication by $z^{-1}$, which is outside the present polynomial
ring unless a localization is added.

The YY generator and the conventional discriminant are
\begin{equation}
X=\frac{1}{1-\psi^5},\qquad \Delta(z)=1-3125z.
\end{equation}
Solving for the same $X$ in $z$ gives
\begin{align}
X&=\frac{3125z}{3125z-1}=-\frac{3125z}{\Delta},\\
X-1&=-\frac1{\Delta},\qquad 1-X=\frac1{\Delta}.
\end{align}
Thus a paper defining its variable as $1/\Delta$ uses $1-X$, not this $X$.
The sign changes under this substitution are real convention changes.
At large complex structure, $z\to0$, $\psi\to\infty$, and $X\to0$.
At the conifold, $\psi^5=1$ and $X$ has a pole.

Differentiating the defining rational function directly,
\begin{align}
\tht X&=\frac{5\psi^5}{(1-\psi^5)^2}
       =5\left(\frac1{1-\psi^5}\right)
           \left(\frac{\psi^5}{1-\psi^5}\right)\\
       &=5X(X-1),\label{eq:thetaX}\\
\theta_z X&=-X(X-1)=X(1-X).
\end{align}
The second factor is $X-1$, since $\psi^5/(1-\psi^5)=X-1$.

\section{Hodge gauge, metric, and connection signs}

The reference holomorphic-form gauge is YY's local residue representative
\begin{equation}
\Omega=5\psi\,\frac{x_5\,\dd x_1\wedge\dd x_2\wedge\dd x_3}
                         {\del p/\del x_4}.
\end{equation}
An overall period normalization is additionally fixed by the Yukawa coefficient
specified below. The engine stores neither the form nor the periods themselves.
Fix a period vector $\Pi$ and symplectic matrix $\Sigma$ with the convention
\begin{equation}
E:=e^{-K}=-i\,\bar\Pi^{,T}\Sigma\Pi,\qquad
G:=G_{\psi\bar\psi}=\del_\psi\bdel K.
\end{equation}
The overall constant in $E$ cancels from all $B_p$ below. Define the
standard holomorphic Christoffel component by
\begin{equation}
\Gamma_{\rm std}=G^{-1}\del_\psi G.
\end{equation}
YY also use a Christoffel symbol with the opposite sign,
$\Gamma_{\rm YY}=-\Gamma_{\rm std}$. We avoid any ambiguity by using
\begin{equation}
A=\tht\log G=\psi\Gamma_{\rm std}
              =-\psi\Gamma_{\rm YY},\qquad
B=\tht\log E=-\psi\del_\psi K.
\label{eq:ABconnections}
\end{equation}
For a section $h$ of $(T^*)^n\otimes\mathcal L^m$, our covariant derivative is
\begin{equation}
D_\psi h=\del_\psi h-n\Gamma_{\rm std}h+mK_\psi h.
\label{eq:cov-original}
\end{equation}
There is one minus Christoffel contribution for each lower holomorphic tensor
index. Since $K_\psi=-B/\psi$,
\begin{equation}
\psi D_\psi h=\tht h-nAh-mBh.
\end{equation}
Equivalently, YY write $+n\Gamma_{\rm YY}$; the two formulas are identical.

Under $\Omega\mapsto f(\psi)\Omega$, one has
\begin{equation}
K\mapsto K-\log f-\log\bar f,\qquad
F_g\mapsto f^{2-2g}F_g.
\end{equation}
The plus sign in $+mK_\psi$ in \eqref{eq:cov-original} is required:
the derivative of $f^m$ cancels the change $-m\del_\psi\log f$
in the connection. For $F_g$ the Hodge power is $m=2-2g$.

\section{Yukawa normalization and its derivative}

Write
\begin{equation}
C_{\psi\psi\psi}=\kappa\frac{\psi^2}{1-\psi^5},\qquad
\kappa=\frac{(2\pi i)^3}{5^3}
\end{equation}
in the YY period normalization. The coefficient $\kappa$ is an overall
period/gauge normalization; the polynomial engine never numerically
evaluates it. The scalar normalization used by the code is
\begin{equation}
C:=\psi^3C_{\psi\psi\psi}
 =\kappa\frac{\psi^5}{1-\psi^5}=\kappa(X-1).
\end{equation}
Consequently,
\begin{align}
\tht\log C_{\psi\psi\psi}
 &=2+5\frac{\psi^5}{1-\psi^5}=5X-3,\\
\tht\log C&=3+(5X-3)=5X,\qquad \tht C=5XC.
\label{eq:thetaC}
\end{align}
The cancellation of the $+3$ from $\psi^3$ against the $-3$ in the tensor
Yukawa derivative is essential. Omitting it changes every normalized
covariant derivative.

For a coordinate-only change, keeping the same Hodge gauge,
\begin{align}
C_{zzz}&=\left(\frac{\dd\psi}{\dd z}\right)^3 C_{\psi\psi\psi}
 =\left(-\frac{\psi}{5z}\right)^3
       \kappa\frac{\psi^2}{1-\psi^5}\\
 &=-\frac{\kappa}{125z^3}(X-1)
  =\frac{\kappa}{125z^3\Delta}.
\end{align}
The two minus signs, one from the cubic Jacobian and one from
$X-1=-1/\Delta$, cancel. The familiar $5/(z^3\Delta)$ is in a differently
normalized Hodge gauge; it must not be inserted here without also transforming
the amplitudes.

\section{The original differential ring and the geometric closure input}

Define, for $p\ge1$,
\begin{equation}
A_p=\frac{\tht^p G}{G},\qquad
B_p=\frac{\tht^p E}{E},\qquad A=A_1,\quad B=B_1.
\end{equation}
Thus $A_2$ is not $A^2$, and $B_2$ is not $B^2$.
The quotient rule gives the exact identities
\begin{equation}
\tht A_p=A_{p+1}-AA_p,\qquad
\tht B_p=B_{p+1}-BB_p.
\label{eq:quotient}
\end{equation}
The negative terms arise from differentiating $G^{-1}$ or $E^{-1}$.

\subsection{Picard--Fuchs closure of \texorpdfstring{$B_4$}{B4}}
The period equation in this gauge is
\begin{equation}
\left[\tht^4-\psi^{-5}
   (\tht-1)(\tht-2)(\tht-3)(\tht-4)\right]\Pi=0.
\end{equation}
The elementary symmetric sums of $1,2,3,4$ are $10,35,50,24$, so
\begin{equation}
\prod_{j=1}^4(\tht-j)
 =\tht^4-10\tht^3+35\tht^2-50\tht+24.
\end{equation}
Multiplying the period equation by $\psi^5$ and rearranging gives
\begin{equation}
(1-\psi^5)\tht^4\Pi
  =10\tht^3\Pi-35\tht^2\Pi+50\tht\Pi-24\Pi.
\end{equation}
The sign of the constant is negative on the right-hand side. Since the
antiholomorphic period factor is untouched by $\tht$, the same linear
identity holds for $E$. Dividing by $E$ therefore yields
\begin{equation}\boxed{
B_4=10XB_3-35XB_2+50XB-24X.}\label{eq:B4}
\end{equation}
As a check, using \eqref{eq:theta-coordinate}, the same operator becomes
\begin{equation}
\theta_z^4-3125z\prod_{j=1}^4\left(\theta_z+\frac j5\right).
\end{equation}
The factor $(-5)^4=625$ cancels in dividing the entire equation; there
is no extra sign because there are four factors.

\subsection{Special-geometry closure of \texorpdfstring{$A_2$}{A2}}
The chosen quintic geometric closure is YY equation (3.22):
\begin{equation}\boxed{\begin{aligned}
A_2={}&-4B_2-2AB-2B+2B^2-2A\\
      &+10XB+5XA-5X-1.
\end{aligned}}\label{eq:A2}
\end{equation}
For clarity, $-4B_2$ and $+2B^2$ have different meanings.
The origin of the holomorphic part can be exhibited by the
special-geometry identity
\begin{equation}
\bdel\left[-4B_2-A_2-2AB-2B+2B^2-2A+10XB+5XA\right]=0.
\end{equation}
Denoting the bracket by $h(\psi)$, solving for $A_2$ gives the first
line and the $XB,XA$ terms in \eqref{eq:A2}, followed by $-h$.
In the adopted YY gauge the holomorphic input is
\begin{equation}
h=5X+1.
\end{equation}
This explains both signs in $-5X-1$. It is a geometric/gauge input,
not a consequence of the Picard--Fuchs equation alone. Using $+5X+1$
in \eqref{eq:A2} would correspond to solving for the bracket with the wrong sign.

An explicit local consistency check is useful. At large complex structure,
\begin{align}
\omega_0&=1+120z+O(z^2),&
t&=-\log z-770z+O(z^2),\\
B&=-600z+O(z^2),& B_2&=3000z+O(z^2),\\
A&=-1-3850z+O(z^2),& A_2&=1+26950z+O(z^2),\\
X&=-3125z+O(z^2).
\end{align}
Here $B=\tht\log\omega_0$ in the holomorphic limit, and
$A=\tht\log(\del_\psi t)$. The coefficient $770$ follows from
$5(120)(H_5-H_1)=5(120)(77/60)=770$.
Substitution in the bracket gives
\begin{align}
h&=1+[-12000-26950-1200+1200+7700+15625]z+O(z^2)\\
 &=1-15625z+O(z^2)=1+5X+O(z^2).
\end{align}
This checks the displayed constants locally. It does not, by itself,
exclude arbitrary additional holomorphic functions; the exact closure
is the specified geometric input.

\section{Change to \texorpdfstring{$u,v_1,v_2,v_3,X$}{u, v1, v2, v3, X} and the complete derivation table}

The precise change of variables is
\begin{align}
u&=B,\\
v_1&=A+1+2B,\\
v_2&=B_2-B(A+1+2B)=B_2-uv_1,\\
v_3&=B_3-B\{B(1+A+2B)-B_2+5X(1+A+2B)-10X\}\\
   &=B_3+uv_2-5uX(v_1-2).
\label{eq:v3definition}
\end{align}
Its inverse is
\begin{equation}\boxed{\begin{aligned}
B&=u,&A&=v_1-1-2u,\\
B_2&=v_2+uv_1,&B_3&=v_3-uv_2+5uX(v_1-2).
\end{aligned}}\label{eq:inverse}
\end{equation}
The factor $5$ in the last equation is present in the implementation.
The constant $-2$ inside $v_1-2$ is responsible for several terms proportional
to $10X$ and $24X$ below; it cannot be dropped as an irrelevant shift.

\subsection{Derivative of \texorpdfstring{$u$}{u}}
Equation \eqref{eq:quotient} gives
\begin{equation}\boxed{\tht u=B_2-B^2=v_2+uv_1-u^2.}\end{equation}

\subsection{Derivative of \texorpdfstring{$v_1$}{v1}}
First differentiate before substituting:
\begin{align}
\tht v_1&=(A_2-A^2)+2(B_2-B^2)\\
 &=-2B_2-2AB-2B-2A-A^2+10XB+5XA-5X-1.
\end{align}
The $+2B^2$ in $A_2$ cancels the $-2B^2$ from $2\tht B$.
Now substitute \eqref{eq:inverse}:
\begin{align}
\tht v_1={}&-2(v_2+uv_1)-2u(v_1-1-2u)-2u\\
 &-2(v_1-1-2u)-(v_1-1-2u)^2\\
 &+10Xu+5X(v_1-1-2u)-5X-1\\
={}&-v_1^2-2v_2-10X+5Xv_1.
\end{align}
All powers of $u$ cancel. The $-10X$ is the sum of $-5X$ from
$5XA$ and the explicit $-5X$ in $A_2$.

\subsection{Derivative of \texorpdfstring{$v_2$}{v2}}
Using $v_2=B_2-uv_1$ gives
\begin{align}
\tht v_2={}&B_3-uB_2-(\tht u)v_1-u\tht v_1\\
={}&v_3-uv_2+5uX(v_1-2)-uv_2-u^2v_1\\
 &-v_1v_2-uv_1^2+u^2v_1\\
 &+uv_1^2+2uv_2+10uX-5uXv_1\\
={}&v_3-v_1v_2.
\end{align}
The $uv_2$ coefficients are $-1-1+2=0$; the $u^2v_1$ and
$uv_1^2$ terms cancel in pairs; and
$5uX(v_1-2)+10uX-5uXv_1=0$.

\subsection{Derivative of \texorpdfstring{$v_3$}{v3}}
To display all cancellations without an opaque large expansion, put
\begin{equation}
L=-v_2+5X(v_1-2),\qquad B_3=v_3+uL.
\end{equation}
From \eqref{eq:B4} and \eqref{eq:quotient},
\begin{equation}
\tht B_3=(10X-u)(v_3+uL)-35X(v_2+uv_1)+50Xu-24X.
\end{equation}
Independently, using the derivatives already established,
\begin{align}
\tht L
 &=-\tht v_2+5(\tht X)(v_1-2)+5X\tht v_1\\
 &=-v_3+v_1v_2+25X(X-1)(v_1-2)\\
 &\hspace{12mm}+5X(-v_1^2-2v_2-10X+5Xv_1)\\
 &=-v_3+(10X-v_1)L-35Xv_1+50X.
\end{align}
Since $v_3=B_3-uL$,
\begin{align}
\tht v_3={}&10Xv_3-35Xv_2-24X-v_2L\\
 &+u\{-v_3+(10X-v_1)L-35Xv_1+50X-\tht L\}\\
={}&10Xv_3-35Xv_2-24X-v_2[-v_2+5X(v_1-2)]\\
={}&v_2^2-24X-25Xv_2-5Xv_1v_2+10Xv_3.
\end{align}
The two $u^2L$ terms canceled before the second line. The remaining
coefficient of $u$ vanishes by the displayed expression for $\tht L$.
Finally, $-35Xv_2+10Xv_2=-25Xv_2$.

\subsection{The closed table and its meaning}
Combining the results,
\begin{equation}\boxed{\begin{aligned}
\tht u&=v_2+uv_1-u^2,\\
\tht v_1&=-v_1^2-2v_2-10X+5Xv_1,\\
\tht v_2&=-v_1v_2+v_3,\\
\tht v_3&=v_2^2-24X-25Xv_2-5Xv_1v_2+10Xv_3,\\
\tht X&=5X(X-1).
\end{aligned}}\label{eq:table}
\end{equation}
These are ordinary derivation identities in the commutative ring
\begin{equation}
\mathcal R=\QQ[u,v_1,v_2,v_3,X].
\end{equation}
In particular, $\QQ[v_1,v_2,v_3,X]$ is closed under $\tht$. Denote the
restriction of $\tht$ to this subring by $\delta$; explicitly,
\begin{align}
\delta={}&(-v_1^2-2v_2-10X+5Xv_1)\del_{v_1}
 +(v_3-v_1v_2)\del_{v_2}\\
 &+(v_2^2-24X-25Xv_2-5Xv_1v_2+10Xv_3)\del_{v_3}
 +5X(X-1)\del_X.
\end{align}
For a general element of $\mathcal R$,
\begin{equation}
\tht=(v_2+uv_1-u^2)\del_u+\delta.
\end{equation}

Assign the filtration weights
\begin{equation}
\wt(u)=\wt(v_1)=\wt(X)=1,\quad \wt(v_2)=2,\quad\wt(v_3)=3.
\end{equation}
Each image in \eqref{eq:table} has weight at most one greater than the
generator. This is a filtration, not a homogeneous grading preserved by
$\tht$: for example, $\tht X=5X^2-5X$ has both weights two and one.
The YY bound is $\wt(P_g)\le3g-3$ and
$\wt(\Pg gn)\le3g-3+n$ for stable vertices.
Indeed, assuming the bound at lower genus, both
$\Pg{g-1}2$ and every product $\Pg h1\Pg{g-h}1$ have weight at most
$3g-4$. Extracting $u^k$ reduces that bound by $k$. Hence the three
gradient components in \eqref{eq:gradient} have weights at most
$3g-4,3g-5,3g-6$, respectively: multiplication of $r_{g,2}$ by $X$
raises its weight by one in the middle component. Integration in variables
of weights $1,2,3$ then gives weight at most $3g-3$ in every case.
The degree restriction on the added function of $X$ is the additional
polynomial/meromorphic ansatz of the YY construction.

\section{Normalized vertices and the covariant derivative: every factor}

For $g\ge2$, define
\begin{equation}
P_g=\Pg g0=C^{g-1}F_g,\qquad
\Pg gn=C^{g-1}\psi^n D_\psi^n F_g.
\end{equation}
The upper parenthesized index is the number of insertions; it is not a power.
The stable low-genus definitions are
\begin{equation}
\Pg 0n=C^{-1}\psi^nD_\psi^{n-3}C_{\psi\psi\psi}\ (n\ge3),\qquad
\Pg 1n=\psi^nD_\psi^{n-1}\del_\psi F_1\ (n\ge1).
\end{equation}
No zero-insertion $P_0$ or $P_1$ is needed by the algorithm. In particular,
$\Pg 03=1$ does not mean that the genus-zero prepotential is one.

Let $H=\Pg gn$, and let $h$ be the corresponding unnormalized rank-$n$
vertex, so $h=C^{1-g}\psi^{-n}H$. By the product rule and \eqref{eq:thetaC},
\begin{align}
\tht h
 &=C^{1-g}\psi^{-n}
       \{\tht H+[5(1-g)X-n]H\}.
\end{align}
The Hodge power is $m=2-2g$. Equation \eqref{eq:cov-original} then gives
\begin{align}
C^{g-1}\psi^{n+1}D_\psi h
 &=\tht H+[5(1-g)X-n-nA-(2-2g)B]H\\
 &=\tht H+[-n(A+1)+(2g-2)B-5(g-1)X]H\\
 &=\tht H+
   \left[-n(v_1-2u)+(2g-2)\left(u-\frac52X\right)\right]H.
\end{align}
Therefore
\begin{equation}\boxed{
\Pg g{n+1}=\mathcal D_{g,n}\Pg gn,
\quad \mathcal D_{g,n}=\tht-n(v_1-2u)
 +(2g-2)\left(u-\frac52X\right).}\label{eq:cov}
\end{equation}
The ledger of contributions is
\begin{center}
\begin{tabular}{ll}
\toprule
Contribution & Origin\\\midrule
$-n$ & derivative of $\psi^{-n}$\\
$-nA$ & $n$ lower covariant tensor indices\\
$+(2g-2)B$ & Hodge power $2-2g$ and $K_\psi=-B/\psi$\\
$-5(g-1)X$ & derivative of $C^{1-g}$\\
$-n(A+1)=-n(v_1-2u)$ & inverse generator change\\
$5/2$ & $5(g-1)=(2g-2)(5/2)$\\\bottomrule
\end{tabular}
\end{center}
For two derivatives, the second operator is $\mathcal D_{g,1}$,
not another $\mathcal D_{g,0}$. For example,
\begin{equation}
\Pg 22=\mathcal D_{2,1}\mathcal D_{2,0}P_2.
\end{equation}
Genus is unchanged during differentiation; tensor rank increases by one.

The genus-zero seed follows by cancellation:
\begin{equation}
\Pg 03=C^{-1}\psi^3C_{\psi\psi\psi}=1.
\end{equation}
As a sign/rank check,
\begin{equation}
\Pg 04=\mathcal D_{0,3}1
 =-3(v_1-2u)-2(u-5X/2)=5X-3v_1+4u.
\end{equation}

\section{Genus one: the seed and the numerical constants}

The Euler characteristic parameter used by YY and by this implementation is
$\chi=-200$, the original quintic's value in this mirror-symmetry convention.
It must not be replaced by the Euler characteristic $+200$ of the resolved
mirror while keeping all other formulas unchanged. The genus-one anomaly is
\begin{equation}
\bdel\del_\psi F_1
 =\frac12 C_{\psi\psi\psi}\bar C_{\bar\psi}^{\psi\psi}
   +\left(1-\frac\chi{24}\right)G.
\end{equation}
Here
\begin{equation}
\bar C_{\bar\psi}^{\psi\psi}
 =e^{2K}G^{-2}\bar C_{\bar\psi\bar\psi\bar\psi}.
\end{equation}
The bar label is antiholomorphic; the raised indices come from inverse
metrics, not from changing the holomorphic Yukawa normalization.

The adopted genus-one holomorphic boundary data give the local derivative
\begin{equation}
\del_\psi F_1=\frac12\del_\psi
 \left[-\log G+\frac{62}{3}K
       +\frac{62}{3}\log\psi-\frac16\log(1-\psi^5)\right].
\label{eq:F1input}
\end{equation}
Terms depending only on $\bar\psi$ have no effect on this formula.
The coefficient $31/3=2-\chi/24$ gives $62/3$ inside the bracket because
of the prefactor $1/2$. Indeed, using
$\bdel\del_\psi\log G=2G-C_{\psi\psi\psi}\bar C_{\bar\psi}^{\psi\psi}$,
the first two terms give
\begin{equation}
-G+\frac12 C_{\psi\psi\psi}\bar C_{\bar\psi}^{\psi\psi}+\frac{31}{3}G
 =\frac12 C_{\psi\psi\psi}\bar C_{\bar\psi}^{\psi\psi}+\frac{28}{3}G,
\end{equation}
where $28/3=1-\chi/24$. The $\log\psi$ and conifold-log terms
are holomorphic boundary inputs, not determined by this mixed derivative.

Since $\tht\log(1-\psi^5)=-5(X-1)$, multiplication of
\eqref{eq:F1input} by $\psi$ gives
\begin{align}
\Pg 11
 &=-\frac12A-\frac{31}{3}B+\frac{31}{3}
      +\frac5{12}(X-1)\\
 &=\frac{31}{3}(1-B)+\frac5{12}(X-1)-\frac12A\\
 &=\boxed{\frac{125}{12}-\frac12v_1-\frac{28}{3}u+\frac5{12}X}.
\label{eq:seed}
\end{align}
The constants are $31/3-5/12+1/2=125/12$; the $u$ coefficient is
$-31/3+1=-28/3$. The $+5/12$ in the conifold term comes from
$\tfrac12(-\tfrac16)(-5)$: both minus signs are needed.

\section{From the geometric BCOV equation to the polynomial anomaly}

For $g\ge2$, the BCOV equation \cite{BCOV} is
\begin{equation}
\bdel F_g=\frac12\bar C_{\bar\psi}^{\psi\psi}
 \left[D_\psi D_\psi F_{g-1}
       +\sum_{h=1}^{g-1}D_\psi F_hD_\psi F_{g-h}\right].
\label{eq:BCOV}
\end{equation}
The $1/2$ is the BCOV symmetry factor. The sum as written is ordered:
$h$ and $g-h$ both occur unless they are equal.

\subsection{Normalization of each term on the right}
For the first term,
\begin{equation}
D_\psi D_\psi F_{g-1}=C^{2-g}\psi^{-2}\Pg{g-1}2.
\end{equation}
For each product,
\begin{align}
D_\psi F_hD_\psi F_{g-h}
 &=C^{1-h}\psi^{-1}\Pg h1\,
   C^{1-(g-h)}\psi^{-1}\Pg{g-h}1\\
 &=C^{2-g}\psi^{-2}\Pg h1\Pg{g-h}1.
\end{align}
The exponents add as $(1-h)+(1-g+h)=2-g$.
Since $C$ is holomorphic, $\bdel(C^{g-1}F_g)=C^{g-1}\bdel F_g$.
Multiplying \eqref{eq:BCOV} by $C^{g-1}$ therefore leaves precisely
$C^{g-1}C^{2-g}\psi^{-2}=C/\psi^2$:
\begin{equation}
\bdel P_g=\frac12\frac C{\psi^2}\bar C_{\bar\psi}^{\psi\psi}H_g,
\qquad
H_g=\Pg{g-1}2+\sum_{h=1}^{g-1}\Pg h1\Pg{g-h}1.
\label{eq:normalizedBCOV}
\end{equation}

\subsection{The minus sign from special geometry}
Set $a=\bdel A$ and $b=\bdel B$. From \eqref{eq:ABconnections},
\begin{equation}
b=-\psi G.
\end{equation}
Special geometry in our Christoffel convention gives
\begin{equation}
\bdel\Gamma_{\rm std}=2G-C_{\psi\psi\psi}\bar C_{\bar\psi}^{\psi\psi}.
\end{equation}
Thus
\begin{align}
a+2b
 &=\psi[2G-C_{\psi\psi\psi}\bar C_{\bar\psi}^{\psi\psi}]-2\psi G\\
 &=-\psi C_{\psi\psi\psi}\bar C_{\bar\psi}^{\psi\psi}
  =-\frac C{\psi^2}\bar C_{\bar\psi}^{\psi\psi}.
\end{align}
Inserting this identity into \eqref{eq:normalizedBCOV},
\begin{equation}\boxed{\bdel P_g=-\frac12(a+2b)H_g.}\label{eq:minus}
\end{equation}
The sign is negative even though \eqref{eq:BCOV} has a positive $1/2$.
It comes from the relation between the antiholomorphic Yukawa and $a+2b$.

As an optional propagator check, the adopted holomorphic propagator choice is
\begin{equation}
S^{\psi\psi}=\frac{2K_\psi+2/\psi-\del_\psi\log G}{C_{\psi\psi\psi}},
\qquad
T^{\psi\psi}:=\frac C{\psi^2}S^{\psi\psi}=2-2B-A=3-v_1.
\end{equation}
The $2/\psi$ term is the propagator-gauge contribution corresponding to
the local holomorphic choice $\mathfrak f=\psi$; it is not a second
copy of $K_\psi$. Since $\bdel S^{\psi\psi}=\bar C_{\bar\psi}^{\psi\psi}$,
one obtains again $(C/\psi^2)\bar C=-\bdel v_1=-(a+2b)$.

\subsection{Antiholomorphic changes of all generators}
The quotient identities and $b=-\psi G$ give
\begin{align}
\bdel B_2&=\bdel(\tht B+B^2)
          =\tht b+2Bb=(A+1+2B)b=v_1b,\\
\bdel B_3&=\bdel(\tht B_2+BB_2)\\
 &=\{\tht v_1+(A+1+B)v_1+B_2\}b\\
 &=\{-v_2+5X(v_1-2)\}b.
\end{align}
In the first line $\tht b=(A+1)b$, since $b=-\psi G$.
The second simplification uses the already derived $\tht v_1$.
Consequently,
\begin{align}
\bdel u&=b,\\
\bdel v_1&=a+2b,\\
\bdel v_2&=\bdel(B_2-uv_1)
           =v_1b-v_1b-u(a+2b)=-u(a+2b),\\
\bdel v_3&=\bdel[B_3+uv_2-5uX(v_1-2)]\\
 &=\{-v_2+5X(v_1-2)\}b+v_2b-u^2(a+2b)\\
 &\hspace{8mm}-5X(v_1-2)b-5uX(a+2b)\\
 &=-u(u+5X)(a+2b),\\
\bdel X&=0.
\end{align}
Every term proportional to $b$ in the intermediate expression for
$\bdel v_3$ cancels; the remaining two terms factor as shown.

\subsection{The two formal equations and the three gradient components}
By the chain rule,
\begin{equation}
\bdel P_g=b\del_uP_g+(a+2b)\mathscr L P_g,
\quad
\mathscr L=\del_{v_1}-u\del_{v_2}-u(u+5X)\del_{v_3}.
\end{equation}
Comparing this with \eqref{eq:minus} in the formal YY anomaly system gives
\begin{equation}\boxed{\del_uP_g=0,\qquad \mathscr L P_g=-\frac12 H_g.}\label{eq:PDE}
\end{equation}
More explicitly, the coefficient of $a$ gives $\mathscr L P_g=-H_g/2$;
the coefficient of $b$ gives $\del_uP_g+2\mathscr L P_g=-H_g$;
subtracting twice the first equation gives $\del_uP_g=0$.

The coefficient comparison is a formal polynomial-ring statement used by YY.
It should not be misread as claiming that two $(0,1)$ differential forms on a
one-complex-dimensional base are linearly independent over all smooth functions.
The implementation adopts the resulting formal differential/anomaly algebra.

Define the source used by the code by
\begin{equation}
R_g=-\frac12 H_g=r_{g,0}+ur_{g,1}+u^2r_{g,2}.
\end{equation}
The absence of higher powers follows from the insertion formulas discussed
below and is checked by the program. Since $P_g$ is independent of $u$,
\begin{align}
\mathscr L P_g
 &=\del_{v_1}P_g
    +u[-\del_{v_2}P_g-5X\del_{v_3}P_g]
    -u^2\del_{v_3}P_g.
\end{align}
Therefore
\begin{equation}\boxed{\begin{aligned}
\del_{v_1}P_g&=r_{g,0},\\
\del_{v_2}P_g&=-r_{g,1}+5Xr_{g,2},\\
\del_{v_3}P_g&=-r_{g,2}.
\end{aligned}}\label{eq:gradient}
\end{equation}
In the second equation the plus sign comes from substituting
$\del_{v_3}P_g=-r_{g,2}$ into
$-\del_{v_2}P_g-5X\del_{v_3}P_g=r_{g,1}$.
Integrating only $r_{g,0}$ would incorrectly leave arbitrary functions of
$v_2,v_3,X$, whereas the full system leaves only a function of $X$.

\section{Exact gradient integration and the ambiguity convention}

Write $G_i=\del_{v_i}P_g$ as obtained from \eqref{eq:gradient}.
The consistency conditions are
\begin{equation}
\del_{v_j}G_i=\del_{v_i}G_j\qquad(1\le i,j\le3).
\end{equation}
They are tested as exact polynomial equalities. A particular primitive is
constructed in the order $v_3,v_2,v_1$:
\begin{align}
Q_3&=\int G_3\,\dd v_3,\\
Q_{23}&=Q_3+\int(G_2-\del_{v_2}Q_3)\,\dd v_2,\\
P_g^{\an}&=Q_{23}+\int(G_1-\del_{v_1}Q_{23})\,\dd v_1.
\end{align}
Every primitive here has zero integration constant in its integration variable.
The mixed-partial identities ensure that each residual has no dependence on
variables already integrated; a final check verifies all three derivatives.
Equivalently, the selected primitive satisfies
\begin{equation}
P_g^{\an}(v_1=v_2=v_3=0,X)=0.
\end{equation}
This is an algebraic choice of representative, not a physical boundary condition.

The complete polynomial is stored as
\begin{equation}
P_g=P_g^{\an}+a_g(X),\qquad
a_g(X)=\sum_{j=0}^{3g-3}a_{g,j}X^j.
\end{equation}
The vector has up to $3g-2$ entries because both endpoint powers, zero and
$3g-3$, are included. This is the full polynomial kernel allowed by the coarse
weight bound; geometric regularity and pole restrictions may reduce the physical
ambiguity subspace.

The conventional holomorphic ambiguity in $F_g$ obeys
\begin{equation}
\Delta F_g=C^{1-g}\Delta a_g(X).
\end{equation}
Thus a vector of coefficients of $a_g$ is not directly a vector of coefficients
of $f_g(\psi)$ in a different meromorphic basis. Moreover, our $P_g^{\an}$ need
not equal $C^{g-1}$ times a particular Feynman-diagram representative. The
reference pure-$X$ vector includes whichever shift is required by our zero-at-
$v=0$ primitive convention.

A more restrictive meromorphic ansatz used in YY for the holomorphic addition
to a Feynman-diagram representative is
\begin{equation}
f_g(\psi)=\sum_{j=1}^{\lfloor(2g-2)/5\rfloor}b_{g,j}\psi^{-5j}
          +\sum_{j=0}^{2g-2}d_{g,j}(1-\psi^5)^{-j}.
\end{equation}
Here the symbols $d_{g,j}$ avoid confusing this basis with $a_{g,j}$.
Because $\psi^{-5}=X/(X-1)$, multiplication by $C^{g-1}$ gives
\begin{align}
C^{g-1}f_g=\kappa^{g-1}\bigg[&
 \sum_{j=1}^{\lfloor(2g-2)/5\rfloor}
       b_{g,j}X^j(X-1)^{g-1-j}\\
 &+\sum_{j=0}^{2g-2}d_{g,j}(X-1)^{g-1}X^j\bigg].
\end{align}
All exponents of $X-1$ in the first sum are nonnegative for $g\ge2$.
The second sum has maximal degree $(g-1)+(2g-2)=3g-3$, explaining
the coarse polynomial degree used by the code. Physical constraints can
therefore select an affine subspace of the full coefficient vector.
Because our primitive differs from the chosen Feynman representative by a
pure-$X$ polynomial, the displayed restrictions apply directly to differences
of allowed amplitudes or to a consistently converted ambiguity, not blindly
to the raw vector returned by our zero-at-$v=0$ convention.

In the API, \code{ambiguity=None} means unresolved. It is distinct from
\code{ambiguity=()}, which explicitly assigns zero to every coefficient.
Continuing to higher genus requires an explicit choice for every lower genus.
The current code propagates rational choices, not symbolic unknowns. Across
several genera the dependence on earlier ambiguity parameters is generally
nonlinear, so a symbolic extension would need a suitable coefficient algebra,
not merely an unresolved linear vector attached to the output.

\section{The genus-two recursion, fully expanded}
\label{sec:g2}

Put
\begin{equation}
p:=\Pg11=\frac{125}{12}-\frac12v_1-\frac{28}{3}u+\frac5{12}X.
\end{equation}
The genus-two source is
\begin{equation}
H_2=\Pg12+p^2,\qquad
\Pg12=\mathcal D_{1,1}p=\tht p-(v_1-2u)p.
\end{equation}
There is one product term because the ordered sum has only $h=1$;
there is no extra factor two at genus two.

\subsection{Differentiate the seed and apply the rank-one connection}
Using the closure table gives
\begin{align}
\tht p={}&\frac12v_1^2-\frac{25}{3}v_2+\frac{35}{12}X
 -\frac52v_1X+\frac{25}{12}X^2
 -\frac{28}{3}uv_1+\frac{28}{3}u^2.
\end{align}
For example, the $v_2$ coefficient is $1-28/3=-25/3$;
the $X$ coefficient is $5-25/12=35/12$.
The connection contribution is
\begin{align}
-(v_1-2u)p={}&-\frac{125}{12}v_1+\frac12v_1^2
 -\frac5{12}v_1X+\frac{125}{6}u\\
 &+\left(\frac{28}{3}-1\right)uv_1
 -\frac{56}{3}u^2+\frac56uX.
\end{align}
Adding yields
\begin{align}
\Pg12={}&v_1^2-\frac{25}{3}v_2-\frac{125}{12}v_1
 +\frac{35}{12}X-\frac{35}{12}v_1X+\frac{25}{12}X^2\\
 &+u\left(\frac{125}{6}-v_1+\frac56X\right)-\frac{28}{3}u^2.
\label{eq:P12}
\end{align}
In particular the $uv_1$ coefficient is $-28/3+28/3-1=-1$.

\subsection{Square the one-point seed}
The square is an ordinary commutative polynomial square:
\begin{align}
p^2={}&\frac{15625}{144}-\frac{125}{12}v_1+\frac14v_1^2
 +\frac{625}{72}X-\frac5{12}v_1X+\frac{25}{144}X^2\\
 &+u\left(-\frac{1750}{9}+\frac{28}{3}v_1-\frac{70}{9}X\right)
 +\frac{784}{9}u^2.
\end{align}
Every cross term carries its factor two. For example,
$2(125/12)(5/12)=625/72$, and
$2(-1/2)(-28/3)=28/3$.

\subsection{Collect the source and apply the negative half}
Adding \eqref{eq:P12} to the square gives
\begin{align}
[u^0]H_2={}&\frac{15625}{144}-\frac{125}{6}v_1+\frac54v_1^2
 -\frac{25}{3}v_2+\frac{835}{72}X
 -\frac{10}{3}v_1X+\frac{325}{144}X^2,\\
[u^1]H_2={}&-\frac{3125}{18}+\frac{25}{3}v_1-\frac{125}{18}X,\\
[u^2]H_2={}&\frac{700}{9}.
\end{align}
For example,
\begin{align}
\frac{35}{12}+\frac{625}{72}&=\frac{835}{72},&
\frac{25}{12}+\frac{25}{144}&=\frac{325}{144},\\
\frac{125}{6}-\frac{1750}{9}&=-\frac{3125}{18},&
-\frac{28}{3}+\frac{784}{9}&=\frac{700}{9}.
\end{align}
Since $R_2=-H_2/2$,
\begin{align}
r_{2,0}={}&-\frac{15625}{288}+\frac{125}{12}v_1-\frac58v_1^2
 +\frac{25}{6}v_2-\frac{835}{144}X
 +\frac53v_1X-\frac{325}{288}X^2,\\
r_{2,1}={}&\frac{3125}{36}-\frac{25}{6}v_1+\frac{125}{36}X,\\
r_{2,2}={}&-\frac{350}{9}.
\end{align}
Applying \eqref{eq:gradient} now gives
\begin{align}
\del_{v_1}P_2={}&r_{2,0},\\
\del_{v_2}P_2={}&-\frac{3125}{36}+\frac{25}{6}v_1-\frac{2375}{12}X,\\
\del_{v_3}P_2={}&\frac{350}{9}.
\end{align}
The $X$ coefficient in the second line is
\begin{equation}
-\frac{125}{36}+5\left(-\frac{350}{9}\right)
 =-\frac{125+7000}{36}=-\frac{2375}{12}.
\end{equation}
The positive $350/9$ in the third line is the second minus sign:
$\del_{v_3}P_2=-r_{2,2}$.

\subsection{Integrate and list all ten anomaly coefficients}
First integrate $G_3$, then $G_2$, then the remaining part of $G_1$:
\begin{align}
Q_3={}&\frac{350}{9}v_3,\\
Q_{23}={}&Q_3-\frac{3125}{36}v_2+\frac{25}{6}v_1v_2
                         -\frac{2375}{12}v_2X,\\
P_2^{\an}={}&Q_{23}-\frac{15625}{288}v_1+\frac{125}{24}v_1^2
 -\frac5{24}v_1^3-\frac{835}{144}v_1X\\
 &+\frac56v_1^2X-\frac{325}{288}v_1X^2.
\end{align}
For the last line, the quadratic term of $G_1$ integrates as
$\int(-5v_1^2/8)\dd v_1=-5v_1^3/24$; the mixed term integrates as
$\int(5v_1X/3)\dd v_1=5v_1^2X/6$. The $25v_2/6$ part of $G_1$
has already been supplied by $\del_{v_1}(25v_1v_2/6)$ and must not be
integrated a second time.

\begin{center}
\begin{tabular}{lll}
\toprule
Monomial & Coefficient & Direct source in gradient integration\\\midrule
$v_1$ & $-15625/288$ & constant term of $G_1$\\
$v_1^2$ & $125/24$ & $(125/12)/2$\\
$v_1^3$ & $-5/24$ & $(-5/8)/3$\\
$v_2$ & $-3125/36$ & constant term of $G_2$\\
$v_1v_2$ & $25/6$ & $v_1$ term of $G_2$\\
$v_3$ & $350/9$ & $G_3$\\
$v_1X$ & $-835/144$ & $X$ term of $G_1$\\
$v_1^2X$ & $5/6$ & $(5/3)/2$\\
$v_2X$ & $-2375/12$ & $X$ term of $G_2$\\
$v_1X^2$ & $-325/288$ & $X^2$ term of $G_1$\\\bottomrule
\end{tabular}
\end{center}

\subsection{The four supplied pure-\texorpdfstring{$X$}{X} coefficients}
The reference choice is
\begin{equation}
a_2(X)=\frac{3125}{144}-\frac{28795}{144}X
       +\frac{205}{144}X^2+\frac{25}{48}X^3.
\label{eq:a2reference}
\end{equation}
These are, in order, the four entries of \code{YY_GENUS2_AMBIGUITY}.
The full polynomial is
\begin{equation}\boxed{\begin{aligned}
P_2={}&\frac{3125}{144}-\frac{15625}{288}v_1
 +\frac{125}{24}v_1^2-\frac5{24}v_1^3
 -\frac{3125}{36}v_2+\frac{25}{6}v_1v_2+\frac{350}{9}v_3\\
 &-\frac{28795}{144}X-\frac{835}{144}v_1X
 +\frac56v_1^2X-\frac{2375}{12}v_2X
 +\frac{205}{144}X^2-\frac{325}{288}v_1X^2+\frac{25}{48}X^3.
\end{aligned}}\label{eq:P2full}
\end{equation}
Equation \eqref{eq:P2full} agrees with YY equation (4.1).
The anomaly equations annihilate every term in \eqref{eq:a2reference}.
Consequently they cannot determine $-28795/144$, $205/144$, or $25/48$,
or the constant $3125/144$, without further information. The program imports
the complete reference vector rather than claiming a boundary-fitting derivation
of those three remaining numbers.

There is a useful independent check of the constant coefficient. In the
A-model convention specified in Section~\ref{sec:amodel}, the genus-two
constant-map value is $F_2^A(0)=\chi/5760=-5/144$. The conversion factor
at large complex structure is $-1/625$, so
\begin{equation}
-\frac1{625}\,a_{2,0}=-\frac5{144}
\quad\Longrightarrow\quad a_{2,0}=\frac{3125}{144}.
\end{equation}
This verifies the imported constant within the stated A-model sign convention.
It is not an implementation of all genus-two boundary constraints.

\section{Why the source is quadratic in \texorpdfstring{$u$}{u} and an all-genus check}

For $g\ge2$ let $k=2g-2$, $p=P_g$, and define
\begin{equation}
\mathcal Q_g=\delta-\frac52kX=\delta-5(g-1)X.
\end{equation}
Here $p$ is independent of $u$, so
\begin{equation}
\Pg g1=\mathcal Q_gp+ku p.
\end{equation}
Applying $\mathcal D_{g,1}$ and using $\tht u=v_2+uv_1-u^2$ gives
\begin{equation}\boxed{\begin{aligned}
\Pg g2={}&(\delta-v_1-\tfrac52kX)\mathcal Q_gp+kv_2p\\
 &+(2k+2)u\,\mathcal Q_gp+k(k+1)u^2p.
\end{aligned}}\label{eq:quadraticvertex}
\end{equation}
For completeness, the linear-$u$ terms before simplification are
\begin{equation}
k\delta p-\frac52k^2Xp+(k+2)\mathcal Q_gp
  =(2k+2)\mathcal Q_gp,
\end{equation}
and the quadratic coefficient is $-k+k(k+2)=k(k+1)$.
The $+kuv_1p$ from $k(\tht u)p$ cancels the $-kuv_1p$ from
the rank-one connection.
The genus-one seed is separately linear in $u$, and \eqref{eq:P12}
is quadratic. Therefore every term of $H_g$ is at most quadratic in $u$.

For an entirely explicit higher-genus coefficient prescription, set
$s=125/12-v_1/2+5X/12$, $c=-28/3$, and, for $h\ge2$,
\begin{equation}
k_h=2h-2,\qquad q_h=\mathcal Q_hP_h,\qquad
t_h=(\delta-v_1-\tfrac52k_hX)q_h+k_hv_2P_h.
\end{equation}
For $g\ge3$, with $m=g-1$, collecting the endpoint terms separately gives
\begin{align}
[u^0]H_g={}&t_m+2s q_m
             +\sum_{h=2}^{g-2}q_hq_{g-h},\\
[u^1]H_g={}&(2k_m+2)q_m+2(sk_mP_m+cq_m)\\
 &+\sum_{h=2}^{g-2}
       (k_hP_hq_{g-h}+k_{g-h}P_{g-h}q_h),\\
[u^2]H_g={}&k_m(k_m+1)P_m+2ck_mP_m\\
 &+\sum_{h=2}^{g-2}k_hk_{g-h}P_hP_{g-h}.
\end{align}
Each interior sum is ordered, and it is empty at $g=3$. The endpoint
factor two is the pair $h=1,m$ and $h=m,1$; it is not present at genus
two, where they are the same single term. Multiplying each of these three
lines by $-1/2$ gives $r_{g,0},r_{g,1},r_{g,2}$, respectively.
These formulas provide a direct route to a future blocked implementation
without first building an expanded polynomial in $u$.

The coefficient $c$ of $u$ in the genus-one seed is
\begin{equation}
c=\frac\chi{24}-1=-\frac{28}{3}.
\end{equation}
For $g=2$, $[u^2]H_2=c+c^2=c(c+1)$, hence
\begin{equation}
\del_{v_3}P_2=\frac12c(c+1)
 =\frac12\left(-\frac{28}{3}\right)\left(-\frac{25}{3}\right)
 =\frac{350}{9}.
\end{equation}
For $g\ge3$, put $a=2g-4$. Equation \eqref{eq:quadraticvertex} and
the two endpoint terms $h=1,g-1$ of the ordered genus sum yield
\begin{align}
[u^2]H_g={}&a(a+1)P_{g-1}+2caP_{g-1}\\
 &+\sum_{h=2}^{g-2}(2h-2)(2g-2h-2)P_hP_{g-h}.
\end{align}
As $\del_{v_3}P_g=[u^2]H_g/2$, this becomes
\begin{equation}\boxed{\begin{aligned}
\del_{v_3}P_g={}&(g-2)\left(2g-5+\frac\chi{12}\right)P_{g-1}\\
 &+2\sum_{h=2}^{g-2}(h-1)(g-h-1)P_hP_{g-h},\qquad g\ge3.
\end{aligned}}\label{eq:allgenusv3}
\end{equation}
The interior sum is still ordered. Its prefactor two arises from
$(2h-2)(2g-2h-2)/2$, not from reducing the ordered sum to a half-sum.
For the first two cases,
\begin{align}
\del_{v_3}P_3&=-\frac{47}{3}P_2,\\
\del_{v_3}P_4&=-\frac{82}{3}P_3+2P_2^2.
\end{align}
The coefficients follow from $\chi/12=-50/3$:
$1(1-50/3)=-47/3$, $2(3-50/3)=-82/3$.
These identities are tested independently of direct residual evaluation.

They can also be derived from YY equation (4.7) for
\begin{equation}
Z=\exp\left(\sum_{g\ge2}\lambda^{2g-2}P_g\right),\qquad
\mathcal E=\lambda\del_\lambda:
\end{equation}
\begin{equation}
\frac2{\lambda^2}\del_{v_3}Z
 =\left[\left(\frac\chi{12}-1\right)\mathcal E
         +\mathcal E^2+\frac\chi{24}\left(\frac\chi{24}-1\right)\right]Z.
\end{equation}
Writing $W=\log Z$, division by $Z$ uses
$\mathcal E^2Z/Z=\mathcal E^2W+(\mathcal EW)^2$.
The constant term gives genus two, the coefficient of $\lambda^{2g-4}$
gives \eqref{eq:allgenusv3}. The formal coupling $\lambda$ counts genus;
it is unrelated to the modulus $z$ and to polynomial weights.

\section{Genus three and the role of supplied genus-two data}

At genus three the source is
\begin{equation}
R_3=-\frac12\left[\Pg22+2\Pg11\Pg21\right].
\end{equation}
The factor two occurs because the ordered terms $h=1$ and $h=2$ are equal.
The genus-two polynomial here includes its chosen $a_2(X)$:
although $a_2$ is annihilated by $\mathscr L$, it is not annihilated by
$\delta$ or by covariant differentiation. For example,
\begin{equation}
\delta a_2(X)=5X(X-1)a_2'(X).
\end{equation}
Consequently genus-three anomaly coefficients depend on the genus-two
boundary vector. Choosing zero genus-two ambiguity instead would produce a
different, mathematically consistent formal continuation.

With \eqref{eq:a2reference}, the selected $P_3^{\an}$ has 55 nonzero terms,
weight at most six, no $u$, and no pure-$X$ terms. It is determined up to
\begin{equation}
a_3(X)=a_{3,0}+a_{3,1}X+\cdots+a_{3,6}X^6.
\end{equation}
The code leaves these seven entries unresolved. In particular,
\begin{equation}
[v_3^2]P_3^{\an}=\frac12\left(-\frac{47}{3}\right)\frac{350}{9}
 =-\frac{8225}{27}.
\end{equation}
The factor $1/2$ here is from integrating a term proportional to $v_3$,
not an additional BCOV symmetry factor.
The complete list of the 55 coefficients is given in Appendix~\ref{app:g3}.

\section{Holomorphic and A-model limits: factors not used by the engine}
\label{sec:amodel}

The recursion produces the normalized B-model object. Recovering the
B-model amplitude in the same gauge is simply
\begin{equation}
F_g=C^{1-g}P_g=\kappa^{1-g}(X-1)^{1-g}P_g.
\end{equation}
For the holomorphic limit based at large complex structure,
\begin{equation}
E^{\hol}\propto\omega_0,\qquad
G^{\hol}\propto\del_\psi t,
\end{equation}
where proportionality factors are antiholomorphic or independent of $\psi$
for the relevant logarithmic derivatives. Thus one may evaluate the generators by
\begin{align}
u^{\hol}&=\frac{\tht\omega_0}{\omega_0},\\
v_1^{\hol}&=\frac{\tht^2t}{\tht t}
                   +2\frac{\tht\omega_0}{\omega_0},\\
v_2^{\hol}&=\frac{\tht^2\omega_0}{\omega_0}-u^{\hol}v_1^{\hol},\\
v_3^{\hol}&=\frac{\tht^3\omega_0}{\omega_0}
                   +u^{\hol}v_2^{\hol}-5u^{\hol}X(v_1^{\hol}-2).
\end{align}
The first term in $v_1^{\hol}$ follows from
$\tht\log(\del_\psi t)=\tht\log(\tht t)-1$; its $-1$
cancels the $+1$ in $v_1=A+1+2B$.

In the YY A-model normalization,
\begin{align}
F_g^A(t)&=\left[\frac{(2\pi i)^3\omega_0^2}{5^7C}\right]^{g-1}
                    P_g^{\hol}\\
 &=\left[\frac{\omega_0^2}{5^4(X-1)}\right]^{g-1}P_g^{\hol}\\
 &=\left[\frac{\omega_0^2(\psi^{-5}-1)}{5^4}\right]^{g-1}P_g^{\hol}.
\end{align}
The simplification uses $C=(2\pi i)^3(X-1)/5^3$, leaving $5^{7-3}=5^4$;
also $1/(X-1)=\psi^{-5}-1$. At large complex structure the factor
in brackets tends to $-1/625$, explaining the sign used in the genus-two
constant-map check.

Equivalently, YY define $W_1^3=\omega_0^2/[625(X-1)]$,
$V_i=W_1^i v_i^{\hol}$, and $Y_1=W_1X$. If a monomial in $P_g$
has weight $w$, its contribution to $F_g^A$ becomes that monomial in
$V_1,V_2,V_3,Y_1$ multiplied by $W_1^{3g-3-w}$.
This is the homogeneous five-generator A-model description; it is different
from the inhomogeneous four-generator polynomial $P_g$.

Neither the period series, the mirror map, nor these frame factors are
numerically evaluated by the prototype. They are subsequent exact-series
operations. The distinction is particularly important when comparing
polynomials quoted in different papers.

\section{Exact sparse representation and the meaning of atomic recipes}

An element is represented by
\begin{equation}
P=\sum_{\alpha\in\mathbb N^5}c_\alpha x^\alpha,\quad
x=(u,v_1,v_2,v_3,X),\quad c_\alpha\in\QQ,
\end{equation}
with the exponent order fixed exactly as displayed. Only nonzero coefficients
are stored. A coefficient is a reduced Python \code{Fraction} with integer
numerator and positive denominator. A negative sign belongs to the numerator.
For example,
\begin{equation}
-\frac{325}{288}v_1X^2
\quad\longleftrightarrow\quad
(0,1,0,0,2)\mapsto\operatorname{Fraction}(-325,288).
\end{equation}
Float input is rejected rather than converted into a rational approximation.
An exponent tuple has five nonnegative integers; negative exponents and
polynomial denominators are not supported by this class.

The elementary algorithms are the mathematical rules
\begin{align}
(P+Q)_\alpha&=p_\alpha+q_\alpha,\\
(PQ)_\gamma&=\sum_{\alpha+\beta=\gamma}p_\alpha q_\beta,\\
\del_{x_i}(c_\alpha x^\alpha)
 &=c_\alpha\alpha_i x^{\alpha-e_i},\\
\int c_\alpha x^\alpha\,\dd x_i
 &=\frac{c_\alpha}{\alpha_i+1}x^{\alpha+e_i},\\
\tht(c_\alpha x^\alpha)
 &=c_\alpha\sum_i\alpha_i x^{\alpha-e_i}\tht x_i.
\end{align}
The last formula has no derivative of $c_\alpha$ because the coefficient
field is $\QQ$. Dependence on the modulus is encoded by the generators,
in particular $X$, rather than by an extra rational function of $z$ in
each coefficient. This is why the earlier general $\QQ(z)$ coefficient
layer is unnecessary for the normalized quintic computation.

All equal monomials are combined and zero coefficients discarded after
each operation. The polynomial layer is expanded and canonical. It is
not an unevaluated symbolic expression tree. Scalar division remains
exact; division by a nonconstant polynomial would require a separate
localized/rational-function layer.

Above it, \code{RecipeArena} stores a directed acyclic graph with nodes
for polynomial leaves, addition, multiplication, and normalized covariant
derivatives labeled by $(g,n)$. Equal node specifications are interned;
commutative binary arguments are ordered by node identifiers. Thus repeated
subexpressions and derivative vertices share nodes. This is structural sharing,
not general factorization or discovery of arbitrary algebraic identities.
\code{expand()} evaluates the needed graph in topological order and caches
exact polynomial values. There is no recursive descent through long graph
chains during this evaluation.

The source sum exploits its symmetry by iterating
$1\le h\le\lfloor g/2\rfloor$ with multiplier two for $h<g-h$
and multiplier one for $h=g-h$. For example,
\begin{equation}
H_4=\Pg32+2\Pg11\Pg31+(\Pg21)^2.
\end{equation}
The central term is not doubled. The overall factor $-1/2$ is applied
after assembling this bracket.

The prototype expands the source before extracting its three $u$
coefficients. Hence its atomic representation delays and shares computations
but does not eliminate expansion. The monomial-derivative cache is bounded
by 8192 entries by default. The recipe arena and expanded node values can
grow until explicitly compacted; \code{compact_recipes()} drops that arena
and its cached vertices while preserving amplitudes. Externally held old
recipe handles can keep an old arena alive.

The polynomial algebra itself is commutative. If one formally adjoins the
operator $\tht$, then $[\tht,M_f]=M_{\tht f}$ for multiplication by $f$.
No arbitrary noncommutative operator algebra is needed by this implementation.

\section{API correspondence, exact checks, and extension boundaries}

\begin{longtable}{>{\raggedright\arraybackslash}p{0.31\textwidth}>{\raggedright\arraybackslash}p{0.62\textwidth}}
\toprule
Code entry point & Mathematical operation or invariant\\\midrule\endhead
\code{Poly} & Canonical sparse element of $\QQ[u,v_1,v_2,v_3,X]$.\\
\code{partial} / \code{integrate} & Polynomial derivative / zero-constant primitive in one named generator.\\
\code{coefficient("u", k)} & Extract $[u^k]P$ and remove the $u^k$ exponent.\\
\code{substitute} & Simultaneous substitution; replacements do not substitute into each other.\\
\code{degree()} & Maximum filtration weight; the zero polynomial uses sentinel $-1$.\\
\code{YYRing.theta} & Leibniz derivation from \eqref{eq:table}.\\
\code{YYRing.theta_z} & $-\tht/5$, not $\del_z$.\\
\code{YYRing.covariant(p,g,n)} & $\mathcal D_{g,n}p$ from \eqref{eq:cov}.\\
\code{BCOV.vertex(g,n)} & Shared recipe for the stable normalized insertion vertex.\\
\code{source_recipe(g)} & Shared recipe for $R_g=-H_g/2$.\\
\code{anomaly_gradient(g)} & All three derivatives in \eqref{eq:gradient}.\\
\code{integrate(g)} & Mixed-partial checks and the particular primitive $P_g^{\an}$.\\
\code{set_ambiguity(g, coeffs)} & Assign rational coefficients of $a_g(X)$; invalidate affected derivative caches.\\
\code{Amplitude.polynomial()} & Return the complete chosen $P_g$, or reject an unresolved vector.\\
\code{residual(g,p)} & $\mathscr Lp-R_g$; exact zero is necessary for the polynomial anomaly equation.\\
\code{save} / \code{load} & Versioned exact JSON amplitudes; restore checks conventions and recomputes anomaly representatives.\\
\bottomrule
\end{longtable}

The residual alone does not test $\del_up=0$ for an arbitrary externally
supplied $p$; the amplitude constructor and the tests impose that condition
separately. Likewise a vanishing residual does not determine $a_g(X)$.
JSON numerator/denominator strings preserve exact integers. Saving uses
atomic file replacement; no executable pickle is used. Restoring validates
the mathematical anomaly representatives against the recorded lower-genus
choices, but it does not certify the physical correctness of those choices.

The 25 tests cover rational arithmetic, rejection of floats, sparse operations,
the Leibniz rule, original $A,B_p$ identities, the filtration, tensor rank,
recipe sharing, gradient integrability, the full published $P_2$, the independent
relations \eqref{eq:allgenusv3}, ambiguity/cache behavior, and checkpoint restoration.
With the reference genus-two vector and explicit zero ambiguity choices at
genera three through seven, a formal run to genus eight gives anomaly-term counts
\begin{equation}
10,\ 55,\ 173,\ 403,\ 830,\ 1502,\ 2548\qquad(g=2,\ldots,8).
\end{equation}
Every residual vanishes exactly. These counts describe the specified formal
run, not a physical boundary-complete genus-eight calculation or a proof of
performance at genus 64.

Useful extensions have clear boundaries: a faster exact coefficient backend;
blocked sparse multiplication; direct use of \eqref{eq:quadraticvertex} to
avoid materializing powers of $u$; modular arithmetic with certified rational
reconstruction; exact boundary-condition solvers; and exact period/mirror-map
series evaluation. A different Calabi--Yau model requires replacing the closure
table, normalization, seed, anomaly vector fields, and ambiguity space together.
Changing $\chi$ alone does not generalize the quintic engine.

\appendix
\section{Coordinate changes of the metric connection}
For comparisons in the same Hodge gauge, define
\begin{equation}
\widetilde A=\theta_z\log G_{z\bar z},\qquad
\widetilde B_p=\frac{\theta_z^pE}{E}.
\end{equation}
The metric transforms as
$G_{\psi\bar\psi}=|\del_\psi z|^2G_{z\bar z}$.
Only the holomorphic Jacobian contributes to $\tht\log G$:
\begin{align}
A&=-5\widetilde A+\tht\log(-5z/\psi)
  =-5\widetilde A-6,\\
B_p&=(-5)^p\widetilde B_p.
\end{align}
The $-6$ consists of $-5$ from $\tht\log z$ and $-1$
from $-\tht\log\psi$. Consequently,
\begin{align}
u&=-5\widetilde B_1,\\
v_1&=-5(\widetilde A+1+2\widetilde B_1),\\
v_2&=25\{\widetilde B_2-
    \widetilde B_1(\widetilde A+1+2\widetilde B_1)\}.
\end{align}
For $v_3$, insert these expressions into
$B_3+uv_2-5uX(v_1-2)$ before making any naming convention for new
$z$-generators. In particular, the constant $-2$ does not scale with $v_1$.
A Hodge-gauge change adds its own shift to $B$; the above formulas describe
a coordinate change alone. This is why replacing $\tht$ by $-5\theta_z$
without transforming the metric connection is insufficient.

\section{Complete genus-three anomaly coefficient table}
\label{app:g3}
The following table is generated directly from the exact implementation with
the genus-two boundary vector \eqref{eq:a2reference}. Its entries specify
$P_3^{\an}=\sum c_m m$, with no pure-$X$ terms. Each rational is exact.
It is not a claim that $a_3(X)=0$ is the physical choice.
\begingroup\renewcommand{\arraystretch}{1.5}\setlength{\extrarowheight}{3pt}
\begin{longtable}{lll}
\toprule
Monomial $m$ & Coefficient $c_m$ & Weight\\
\midrule\endhead
$v_2$ & $\dfrac{390625}{864}$ & 2 \\
$v_1 X$ & $-\dfrac{2882375}{576}$ & 2 \\
$v_3$ & $-\dfrac{146875}{432}$ & 3 \\
$v_2 X$ & $-\dfrac{808975}{144}$ & 3 \\
$v_1 X^{2}$ & $-\dfrac{497125}{576}$ & 3 \\
$v_1 v_2$ & $-\dfrac{1990625}{1728}$ & 3 \\
$v_1^{2} X$ & $\dfrac{2917175}{3456}$ & 3 \\
$v_1^{3}$ & $-\dfrac{1953125}{10368}$ & 3 \\
$v_3 X$ & $\dfrac{1353365}{432}$ & 4 \\
$v_2 X^{2}$ & $-\dfrac{1736425}{108}$ & 4 \\
$v_2^{2}$ & $-\dfrac{359375}{288}$ & 4 \\
$v_1 X^{3}$ & $-\dfrac{570875}{1728}$ & 4 \\
$v_1 v_3$ & $\dfrac{734375}{864}$ & 4 \\
$v_1 v_2 X$ & $-\dfrac{6462385}{1728}$ & 4 \\
$v_1^{2} X^{2}$ & $\dfrac{1362425}{3456}$ & 4 \\
$v_1^{2} v_2$ & $-\dfrac{78125}{432}$ & 4 \\
$v_1^{3} X$ & $-\dfrac{801575}{3456}$ & 4 \\
$v_1^{4}$ & $\dfrac{15625}{288}$ & 4 \\
$v_3 X^{2}$ & $-\dfrac{9635}{432}$ & 5 \\
$v_2 X^{3}$ & $\dfrac{1525}{27}$ & 5 \\
$v_2 v_3$ & $\dfrac{146875}{108}$ & 5 \\
$v_2^{2} X$ & $-\dfrac{2951365}{432}$ & 5 \\
$v_1 X^{4}$ & $-\dfrac{200375}{1728}$ & 5 \\
$v_1 v_3 X$ & $\dfrac{39245}{432}$ & 5 \\
$v_1 v_2 X^{2}$ & $-\dfrac{460135}{1728}$ & 5 \\
$v_1 v_2^{2}$ & $-\dfrac{26375}{216}$ & 5 \\
$v_1^{2} X^{3}$ & $\dfrac{560425}{3456}$ & 5 \\
$v_1^{2} v_3$ & $-\dfrac{5875}{72}$ & 5 \\
$v_1^{2} v_2 X$ & $\dfrac{95315}{432}$ & 5 \\
$v_1^{3} X^{2}$ & $-\dfrac{394925}{3456}$ & 5 \\
$v_1^{3} v_2$ & $\dfrac{8125}{144}$ & 5 \\
$v_1^{4} X$ & $\dfrac{3995}{96}$ & 5 \\
$v_1^{5}$ & $-\dfrac{625}{96}$ & 5 \\
$v_3 X^{3}$ & $-\dfrac{1175}{144}$ & 6 \\
$v_3^{2}$ & $-\dfrac{8225}{27}$ & 6 \\
$v_2 X^{4}$ & $\dfrac{2375}{288}$ & 6 \\
$v_2 v_3 X$ & $\dfrac{111625}{36}$ & 6 \\
$v_2^{2} X^{2}$ & $-\dfrac{6774925}{864}$ & 6 \\
$v_2^{3}$ & $-\dfrac{5225}{81}$ & 6 \\
$v_1 X^{5}$ & $-\dfrac{11875}{288}$ & 6 \\
$v_1 v_3 X^{2}$ & $\dfrac{15275}{864}$ & 6 \\
$v_1 v_2 X^{3}$ & $-\dfrac{31775}{1728}$ & 6 \\
$v_1 v_2 v_3$ & $-\dfrac{1175}{18}$ & 6 \\
$v_1 v_2^{2} X$ & $\dfrac{56905}{216}$ & 6 \\
$v_1^{2} X^{4}$ & $\dfrac{163375}{3456}$ & 6 \\
$v_1^{2} v_3 X$ & $-\dfrac{235}{18}$ & 6 \\
$v_1^{2} v_2 X^{2}$ & $\dfrac{775}{108}$ & 6 \\
$v_1^{2} v_2^{2}$ & $\dfrac{1465}{72}$ & 6 \\
$v_1^{3} X^{3}$ & $-\dfrac{317375}{10368}$ & 6 \\
$v_1^{3} v_3$ & $\dfrac{235}{72}$ & 6 \\
$v_1^{3} v_2 X$ & $\dfrac{925}{144}$ & 6 \\
$v_1^{4} X^{2}$ & $\dfrac{875}{72}$ & 6 \\
$v_1^{4} v_2$ & $-\dfrac{15}{4}$ & 6 \\
$v_1^{5} X$ & $-\dfrac{275}{96}$ & 6 \\
$v_1^{6}$ & $\dfrac{5}{16}$ & 6 \\
\bottomrule
\end{longtable}
\endgroup

\section{Reproducibility and source identification}
The accompanying audit script \code{verify_bcov_notes.py} checks the expanded
genus-one second vertex, the three genus-two source coefficients, the complete
genus-two reference polynomial, the general quadratic-insertion identity,
the genus-three/four $v_3$ relations, and the frame-factor simplifications.
The SHA-256 digest below identifies the original 0.1.0 core implementation
that was audited on 16 October 2026. The additional frame and enumerative
checks in Section~\ref{sec:ledger2026} use separate v0.2 modules:
\begin{center}\small\texttt{bd9d0415c8cc0a0fef5b8dd8e085f41538d6115cdb978f3ca569dccb86768313}\end{center}

The formulas are based on YY equations (2.4), (2.12), (2.16), (2.19)--(2.20),
(3.1)--(3.8), (3.16), (3.22)--(3.30), (3.31)--(3.33), (4.1), and (4.7).
The BCOV equation is geometric input; the coefficient arithmetic and the
implementation-specific primitive/ambiguity split are given explicitly here.

\begin{thebibliography}{9}
\bibitem{YY} S.~Yamaguchi and S.-T.~Yau,
\emph{Topological String Partition Functions as Polynomials},
JHEP \textbf{07} (2004) 047,
\href{https://arxiv.org/abs/hep-th/0406078}{arXiv:hep-th/0406078}, version 2.
\bibitem{HKQ} M.~Huang, A.~Klemm and S.~Quackenbush,
\emph{Topological String Theory on Compact Calabi--Yau: Modularity and Boundary Conditions},
\href{https://arxiv.org/abs/hep-th/0612125}{arXiv:hep-th/0612125}.
\bibitem{BCOV} M.~Bershadsky, S.~Cecotti, H.~Ooguri and C.~Vafa,
\emph{Kodaira--Spencer Theory of Gravity and Exact Results for Quantum String Amplitudes},
Commun. Math. Phys. \textbf{165} (1994) 311--428,
\href{https://arxiv.org/abs/hep-th/9309140}{arXiv:hep-th/9309140}.
\end{thebibliography}
\end{document}
